\documentclass[11pt,a4paper]{article}

\usepackage[margin=0.8in]{geometry}
\usepackage{setspace}
\usepackage{enumitem}
\usepackage{fancyhdr}
\usepackage{xcolor}
\usepackage{titlesec}
\usepackage{amssymb}

\definecolor{darkblue}{RGB}{35,65,105}

\setstretch{1.05}

\pagestyle{fancy}
\fancyhf{}
\lhead{\textbf{Yashjot Kaur}}
\rhead{\textcolor{darkblue}{Week 1 Journal}}
\cfoot{\thepage}

\renewcommand{\headrulewidth}{0.5pt}

\titleformat{\section}
{\large\bfseries\color{darkblue}}
{}{0em}{}

\begin{document}

% ---------- TITLE ----------
\begin{center}
    {\LARGE \textbf{\textcolor{darkblue}{WEEK 1 JOURNAL}}}\\[5pt]
    
    {\large \textbf{AI-based Real-Time Sensitive Data Detection and Masking}}\\[7pt]
    
    \textbf{Name: Yashjot Kaur \quad | \quad Roll No.: 1024170452}
\end{center}

\vspace{5pt}
\hrule
\vspace{8pt}

% ---------- OBJECTIVE ----------
\section*{1. Objective}

The main objective of Week 1 was to understand the problem statement,
requirements, and overall idea of the project. I focused on understanding
why sensitive information needs to be detected and protected automatically.

% ---------- WORK DONE ----------
\section*{2. Work Done}

During the first week, I studied the concept of sensitive data and
Personally Identifiable Information (PII). I identified common examples
such as names, phone numbers, email addresses, addresses, and identification
details.

I also explored how Artificial Intelligence and Natural Language Processing
can help in identifying sensitive information from text and how the
identified information can later be masked.

% ---------- PROJECT UNDERSTANDING ----------
\section*{3. Project Understanding}

The basic idea of the proposed system is to identify sensitive information
from the given text and protect it before the information is displayed or
shared.

\vspace{5pt}

\begin{center}

\fbox{\textbf{Raw Text}}
\hspace{8pt}
$\longrightarrow$
\hspace{8pt}
\fbox{\textbf{Detect PII}}
\hspace{8pt}
$\longrightarrow$
\hspace{8pt}
\fbox{\textbf{Mask Data}}
\hspace{8pt}
$\longrightarrow$
\hspace{8pt}
\fbox{\textbf{Safe Output}}

\end{center}

\vspace{5pt}

For example, if the input contains an email address, the system identifies
it as sensitive information and replaces it with a masked value such as
\texttt{[EMAIL MASKED]}.

% ---------- CONTRIBUTION ----------
\section*{4. My Contribution}

My contribution during Week 1 was mainly focused on understanding the
project requirements and identifying the different types of sensitive
information that the system should detect.

I also helped in discussing the overall workflow and organising the
initial project structure for further development.

% ---------- LEARNING ----------
\section*{5. Learning Outcomes}

This week helped me understand the importance of data privacy and the role
of AI in protecting sensitive information. I learned the basic concepts
of PII detection, Natural Language Processing, and data masking.

I also gained a clearer understanding of how the different stages of the
project will work together.

% ---------- NEXT WEEK ----------
\section*{6. Plan for Week 2}

In Week 2, I planned to study suitable detection techniques, prepare
sample inputs, explore NLP-based approaches, and understand how different
types of sensitive information can be detected.

\vfill

\begin{center}
    \textcolor{darkblue}{\rule{0.35\textwidth}{0.5pt}}\\[3pt]
    \textit{\small End of Week 1}
\end{center}

\end{document}